\scanpage{85}
\mspage{42}
moreover $G_m(S)$ has no infinitely divisible element\deleted{\unsure{s}}, hence
\[
\unsure{(22)|(21)}\quad G_m(S)\longrightarrow G_m(S)^\wedge\simeq H^1(S,T)
\]
\margin{(``arithmetic'' case)}
is injective. Thus, when $\gamma_i=0$ i.e. when (13) admits ``\unsure{geometric|geometric}' splittings \inserted{(\unsure{for} it suffices for this that a section exists at all)} these form a torsor under $\deleted{\unsure{H^1(S,T)}\unsure{=|\simeq}}G_m(S)^{\unsure{\wedge}}$, which is identified with a subtorsor of the $(H^1(S,T)=G_m(S)^\wedge)$-torsor of all splittings of (13). For the ``profinite description of absolute algebraic geometry over $\mathbb Q$'' to be complete, one would also have to characterize (in terms of \inserted{this} profinite description) this remarkable subset.

Finally, I would also like to understand how a section of \unsure{an extension|extensions} of type (1) can ``specialize'' to a section of type (2), in the case of \unsure{fields with a} valuation. For this, I take up again
\margin{84}
